\documentclass[11pt,letterpaper]{article}
\usepackage[margin=0.55in]{geometry}
\usepackage{fontspec}
\setmainfont[Path=/usr/local/texlive/2026/texmf-dist/fonts/opentype/public/lm/,BoldFont=lmroman10-bold.otf,ItalicFont=lmroman10-italic.otf]{lmroman10-regular.otf}
\newfontfamily\chinesefont[Path=/System/Library/Fonts/Supplemental/]{Songti.ttc}
\usepackage[colorlinks=true,urlcolor=blue]{hyperref}
\usepackage{tabularx,enumitem,titlesec}
\usepackage{needspace}
\pagestyle{plain}
\emergencystretch=2em
\raggedbottom
\brokenpenalty=10000
\setlength{\parindent}{0pt}
\setlength{\parskip}{1pt}
\titleformat{\section}{\large\bfseries}{}{0pt}{}[\vspace{2pt}\titlerule]
\titlespacing*{\section}{0pt}{4pt}{2pt}
\setlist[itemize]{leftmargin=14pt,itemsep=1pt,topsep=2pt,parsep=0pt}
\newcommand{\entry}[3]{\begin{tabularx}{\linewidth}{@{}p{1.24in}X@{}}\parbox[t]{1.24in}{\textbf{#1}}&\textbf{#2}\newline #3\end{tabularx}\par\vspace{2pt}}
\newcommand{\paper}[1]{\textbf{#1}\par\nopagebreak}
\begin{document}
\fontsize{10.5}{11.5}\selectfont
\begin{center}{\LARGE\bfseries Sunny Liang Tan {\chinesefont (谭靓)}}\end{center}
\begin{tabularx}{\linewidth}{@{}l l X r@{}}
\textbf{Placement Director:}&Lee Ohanian&\href{mailto:ohanian@econ.ucla.edu}{ohanian@econ.ucla.edu}&(310) 825-0979\\[2pt]
\textbf{Graduate Administrator:}&Chiara Paz&\href{mailto:chiara@econ.ucla.edu}{chiara@econ.ucla.edu}&(310) 206-1413
\end{tabularx}
\section*{Contact Information}
\begin{tabularx}{\linewidth}{@{}X X@{}}
\textbf{Email:} \href{mailto:tanliang@ucla.edu}{tanliang@ucla.edu}&\textbf{Phone:} \href{tel:+14244029515}{(+1) 424-402-9515}\\[4pt]
\textbf{Address:} UCLA Department of Economics\newline 315 Portola Plaza\newline Los Angeles, CA 90095&\textbf{Website:} \href{https://sunnyliangtan.github.io}{sunnyliangtan.github.io}
\end{tabularx}
\section*{References}
\begin{tabularx}{\linewidth}{@{}X X@{}}
\textbf{Lee Ohanian} (chair)&\textbf{David Baqaee}\\
Distinguished Professor of Economics&Professor of Economics\\
Department of Economics&Department of Economics\\
University of California, Los Angeles&University of California, Los Angeles\\
\href{mailto:ohanian@econ.ucla.edu}{ohanian@econ.ucla.edu}&\href{mailto:baqaee@econ.ucla.edu}{baqaee@econ.ucla.edu}\\[9pt]
\textbf{David Autor}&\textbf{Jon Vogel}\\
Daniel (1972) and Gail Rubinfeld Professor&Professor of Economics\\
Department of Economics&Department of Economics\\
Massachusetts Institute of Technology&University of California, Los Angeles\\
\href{mailto:dautor@mit.edu}{dautor@mit.edu}&\href{mailto:jvogel@econ.ucla.edu}{jvogel@econ.ucla.edu}
\end{tabularx}
\section*{Education}
\begin{tabularx}{\linewidth}{@{}p{1.24in}X r@{}}
\textbf{Oct 2021 -}&\textbf{University of California, Los Angeles (UCLA)}\newline Ph.D. in Economics\newline Expected Graduation: Jun 2027&Los Angeles, CA
\end{tabularx}\par\vspace{2pt}
\begin{tabularx}{\linewidth}{@{}p{1.24in}X r@{}}
\parbox[t]{1.24in}{\textbf{Sep 2024 -\\May 2025}}&\textbf{Massachusetts Institute of Technology (MIT)}\newline Visiting Scholar&Cambridge, MA
\end{tabularx}\par\vspace{2pt}
\begin{tabularx}{\linewidth}{@{}p{1.24in}X r@{}}
\parbox[t]{1.24in}{\textbf{Sep 2019 -\\May 2021}}&\textbf{Duke University}\newline M.A. in Economics&Durham, NC
\end{tabularx}\par\vspace{2pt}
\begin{tabularx}{\linewidth}{@{}p{1.24in}X r@{}}
\parbox[t]{1.24in}{\textbf{Sep 2015 -\\May 2019}}&\textbf{The University of Hong Kong}\newline B.A. in Economics \& Finance; Minor in Mathematics&Hong Kong
\end{tabularx}\par\vspace{2pt}
\section*{Research Fields}
Primary: Macroeconomics, International Economics\\
Secondary: Development Economics, Industrial Organization\\
Research interests: Industrial Policy, Innovation Policy, Production Networks, GE Sufficient Statistics
\section*{Job Market Paper}
\paper{Government Venture Capital, Production Networks, and Aggregate Productivity in China}
Abstract: Industrial policy can reshape production well beyond the firms it directly supports, and assessing its aggregate productivity implications therefore requires measuring how supply-chain spillovers amplify or counteract recipient gains. This paper provides the first joint analysis of these direct, upstream, and downstream contributions in the setting of government venture capital in China. Linking investment records to firm accounts, patents, and supplier--buyer transactions, we identify expansion among financed firms and their existing trading partners. Suppliers expand sales and capital, while buyers expand sales and payroll. We combine these responses with baseline allocation statistics in a production-network framework that values both the resources absorbed by expanding firms and those released elsewhere. In the benchmark accounting with fixed industry factor supplies, supplier spillovers reinforce the recipient contribution, while the larger buyer contribution counteracts and more than offsets the positive recipient and supplier contributions at the point estimates. The analysis shows that supply-chain spillovers can change both the scale and direction of the aggregate productivity contribution suggested by recipient outcomes alone, highlighting the value of considering production networks when evaluating industrial policy.
\Needspace{7\baselineskip}
\section*{Working Paper}
\textbf{Exchange-Rate Movements and U.S. Corporations} with Ivo Welch and Yuqing Zhou\par\nopagebreak

Abstract: This paper revisits the exchange-rate exposure puzzle: why prior research finds weak relationships between currency movements and U.S. firms' stock returns despite substantial international trade. Linking transaction-level U.S. Census export records to financial statements and stock returns for publicly traded U.S. exporters from 1992 to 2016, we construct firm-specific exchange-rate exposure measures that reflect each firm's export destinations. Dollar depreciation is associated with economically significant increases in exports, sales, profitability, and stock returns, with export, sales, and stock-return sensitivities generally larger among larger and more export-oriented firms. In asset- or export-weighted specifications, a 1\% depreciation is associated with approximately 0.4-0.6\% higher exports and 0.6-0.8\% higher sales. The stronger sales response is consistent with currency movements affecting domestic competitiveness as well as foreign demand, while the profitability results indicate that financial and operational hedging do not fully offset exposure. Real operating sensitivities increase over time. Stock prices respond contemporaneously, with no detectable delayed response around subsequent earnings announcements. By measuring exposure using firms' actual trading relationships and tracing its implications across operating outcomes and equity valuations, the paper contributes new evidence to the exchange-rate exposure literature and shows that substantial corporate currency exposure can coexist with timely incorporation of information into stock prices.
\section*{Work in Progress}
\textbf{Industrial Policy and Scale Complementarities in EV Battery Supply Chains} with Junjie Xia and Binkai Chen\par\nopagebreak

Abstract: This project studies whether scale economies at successive production stages amplify industrial subsidies through mutually reinforcing demand and cost linkages in EV battery supply chains. Support for battery producers can increase demand for materials, allowing suppliers to lower unit costs; when passed downstream, these savings can induce further expansion. Using linked income-tax, VAT, and firm-to-firm transaction data, the analysis traces policy exposure through existing supplier relationships and examines suppliers' scale, costs, investment, and sales to other buyers. It distinguishes scale economies from changes in demand, markups, product composition, and supplier selection, and will discipline a production-network model to compare upstream, downstream, and joint support under a common fiscal budget.
\par\textbf{Industrial Policy and Induced Innovation in Semiconductor Supply Chains} with Meryvn Ho\par\nopagebreak
Abstract: This project studies whether support for downstream chipmakers induces technological upgrading among semiconductor equipment suppliers. More advanced chip production raises equipment requirements, potentially making costly supplier innovation worthwhile, while repeated testing and customer feedback may generate learning that benefits subsequent buyers. Combining government venture capital investments, administrative firm data, supplier--buyer transactions, and fabrication-capacity and process-technology data, the analysis distinguishes customer capacity expansion from technological advances and examines how exposure duration and intensity shape supplier upgrading. It investigates whether new capabilities extend to other buyers and separates upgrading within incumbent suppliers from substitution toward suppliers that already possess advanced capabilities, tracing transferable learning beyond directly financed firms.
\par\textbf{Are Banks Discriminatory? Evidence from U.S. Residential Lending} with Jon Vogel\par\nopagebreak
Abstract: This paper investigates discrimination in U.S. mortgage lending by examining the relationship between credit allocation and subsequent loan performance. Differences in lending across racial groups can reflect either differences in credit risk or unequal lending standards, making lending outcomes alone insufficient to distinguish between these explanations. We combine mortgage lending records, consent-order information, and bank regulatory data to examine whether changes in lending to minority borrowers are accompanied by changes in default and charge-off rates. The central hypothesis is that, if banks apply excessively restrictive standards to minority applicants, expanding lending to these borrowers need not produce the deterioration in loan performance implied by risk-based explanations for restricted credit access. The analysis uses enforcement episodes to study adjustments in lending and subsequent losses, accounting for the timing of loss realization and changes in portfolio composition. By bringing lending decisions and realized credit losses into a common framework, the project aims to assess whether observed disparities are consistent with differential lending standards and to clarify the implications of discrimination for credit allocation.
\Needspace{8\baselineskip}
\section*{Teaching Experience}
\textbf{University of California, Los Angeles (UCLA)}\\
Teaching Assistant, Principles of Economics (Winter 2024, Spring 2026)\\
Teaching Assistant, Microeconomic Theory (ECON 101; Fall 2022, Winter 2023, Fall 2023)\\
Teaching Assistant, Microeconomic Theory (ECON 11; Winter 2026)\\
Teaching Assistant, Ph.D. International Macroeconomics (Spring 2024)\\
Teaching Assistant, Competitive Strategy (elective; Spring 2023)

\Needspace{4\baselineskip}
\textbf{Duke University}\\
Teaching Assistant, Ph.D. Econometrics II (Fall 2020)\\
Teaching Assistant, Econometrics and Data Science (Spring 2021)
\Needspace{10\baselineskip}
\section*{Fellowships and Awards}
\begin{tabularx}{\linewidth}{@{}p{1.24in}X@{}}
\textbf{2021--2027}&UCLA Graduate Fellowship\\[2pt]
\textbf{2020}&M.A. Merit Scholar Award, Duke University\\[2pt]
\textbf{2018}&HKU Centennial Scholarship\\[2pt]
\textbf{2015--2016,}\newline\textbf{2017--2018}&Dean's Honors List\\[6pt]
\textbf{2016}&Hong Kong SAR Government Reaching Out Award
\end{tabularx}\par\vspace{2pt}
\Needspace{8\baselineskip}
\section*{Presentations and Professional Service}
\begin{itemize}
\item Presentation: South Cal International Macro Finance Conference, 2026.
\item Presentation: West Coast Micro-to-Macro Development Workshop, 2026.
\item Presentation: MIT Trade Lunch, 2025.
\item Organizing Committee, Southern California Graduate Conference in Applied Economics, 2023.
\item Initiator, UCLA Women in Economics, 2023-2025.
\end{itemize}
\section*{Research Experience}
\begin{tabularx}{\linewidth}{@{}p{1.24in}X r@{}}
\parbox[t]{1.24in}{\textbf{Jan 2023 -\\Feb 2023}}&\textbf{University of California, Los Angeles (UCLA)}\newline Research Assistant to Jon Vogel\newline Theoretical Analysis of Labor-Market Equilibrium&Los Angeles, CA
\end{tabularx}\par\vspace{2pt}
\begin{tabularx}{\linewidth}{@{}p{1.24in}X r@{}}
\parbox[t]{1.24in}{\textbf{Sep 2021 -\\Nov 2021}}&\textbf{Duke University}\newline Research Assistant to David Berger\newline Literature Reviews and Data Construction on Wages and Taxes&Durham, NC
\end{tabularx}\par\vspace{2pt}
\begin{tabularx}{\linewidth}{@{}p{1.24in}X r@{}}
\parbox[t]{1.24in}{\textbf{Jan 2021 -\\Jun 2021}}&\textbf{Duke University}\newline Research Assistant to Matthew Masten\newline Econometric Proof Summaries and Literature Reviews&Durham, NC
\end{tabularx}\par\vspace{2pt}
\begin{tabularx}{\linewidth}{@{}p{1.24in}X r@{}}
\parbox[t]{1.24in}{\textbf{Nov 2018 -\\Jul 2022}}&\textbf{Massachusetts Institute of Technology (MIT)}\newline Research Assistant to David Autor\newline Occupational Data Digitalization, Classification, and Econometric Analysis&Cambridge, MA
\end{tabularx}\par\vspace{2pt}
\section*{Other Research Experience}
\begin{tabularx}{\linewidth}{@{}p{1.24in}X r@{}}
\parbox[t]{1.24in}{\textbf{Jan 2026 -\\May 2026}}&\textbf{SoftBank China Venture Capital}\newline Data Collaboration\newline China Artificial Intelligence and Energy&Hybrid
\end{tabularx}\par\vspace{2pt}
\begin{tabularx}{\linewidth}{@{}p{1.24in}X r@{}}
\textbf{2023}&\textbf{Price Theory Summer Camp}&Chicago
\end{tabularx}\par\vspace{2pt}
\begin{tabularx}{\linewidth}{@{}p{1.24in}X r@{}}
\parbox[t]{1.24in}{\textbf{Jun 2017 -\\Jul 2017}}&\textbf{Bank of China (Hong Kong)}\newline Summer Analyst\newline Bond Connect Due Diligence / Regulatory Materials&Hong Kong
\end{tabularx}\par\vspace{2pt}
\begin{tabularx}{\linewidth}{@{}p{1.24in}X r@{}}
\parbox[t]{1.24in}{\textbf{Dec 2016 -\\Jan 2017}}&\textbf{Muxin Asset Management (Shanghai)}\newline Quantitative Research Intern\newline A-Share Alpha Strategies / Foreign Exchange Research&Shanghai
\end{tabularx}\par\vspace{2pt}
\section*{Languages}
Mandarin (native); English (fluent); Cantonese and Spanish (conversational)
\par\vspace{2pt}{\small\hfill Last updated: November 2026\par}
\end{document}
